\section{神经网络的实现}

\subsection{用 Python 实现两层网络}

一个完整的两层神经网络可以用不到 20 行 Python 代码实现：

\begin{figure}[H]
\centering
\includegraphics[width=0.95\textwidth]{slides-images/slide-033.jpg}
\caption{两层神经网络的 Python 实现：前向传播、损失计算、梯度计算和参数更新\protect\footnotemark}
\end{figure}
\footnotetext{Slides 第34页。}

\begin{lstlisting}[caption={两层神经网络的核心代码}]
# 定义网络维度
N, D_in, H, D_out = 64, 1000, 100, 10
X = np.random.randn(N, D_in)
Y = np.random.randn(N, D_out)
W1 = np.random.randn(D_in, H)
W2 = np.random.randn(H, D_out)

for t in range(2000):
    # 前向传播
    h = X.dot(W1)             # 第一层线性变换
    h_relu = np.maximum(h, 0) # ReLU 激活
    y_pred = h_relu.dot(W2)   # 第二层线性变换

    # 计算损失
    loss = np.square(y_pred - Y).sum()

    # 反向传播（手动计算梯度）
    grad_y_pred = 2.0 * (y_pred - Y)
    grad_W2 = h_relu.T.dot(grad_y_pred)
    grad_h_relu = grad_y_pred.dot(W2.T)
    grad_h = grad_h_relu.copy()
    grad_h[h < 0] = 0
    grad_W1 = X.T.dot(grad_h)

    # 参数更新
    W1 -= learning_rate * grad_W1
    W2 -= learning_rate * grad_W2
\end{lstlisting}

代码中最关键的部分是\textbf{反向传播}：计算损失对 $W_1$ 和 $W_2$ 的解析梯度。这正是本节课的重点。

\subsection{网络容量与正则化}

\begin{figure}[H]
\centering
\includegraphics[width=0.95\textwidth]{slides-images/slide-035.jpg}
\caption{不同隐藏层大小的分类结果：更多神经元 = 更复杂的决策边界 = 更强的拟合能力\protect\footnotemark}
\end{figure}
\footnotetext{Slides 第36页。}

增加隐藏层神经元的数量可以学到更复杂的决策边界。但过大的网络容量会导致过拟合。

\begin{importantbox}{不要用网络大小作为正则化手段}
正确的做法是：选择一个\textbf{略大于需要}的网络，然后通过调节正则化强度 $\lambda$ 来控制过拟合。原因：
\begin{enumerate}
    \item 训练大型网络成本高，如果发现网络太小还要重新训练
    \item 正则化参数 $\lambda$ 可以快速调整
    \item 先让网络有足够容量``记住''数据，再用正则化让它``忘记''不重要的细节
\end{enumerate}
\end{importantbox}

\begin{figure}[H]
\centering
\includegraphics[width=0.95\textwidth]{slides-images/slide-037.jpg}
\caption{正则化强度对决策边界的影响：$\lambda$ 增大使边界更平滑，但过大会导致欠拟合\protect\footnotemark}
\end{figure}
\footnotetext{Slides 第38页。}

\subsection{本章小结}

实现一个两层神经网络只需要矩阵乘法、ReLU 和损失计算。关键挑战是梯度的计算。网络大小应略大于需要，通过正则化控制过拟合而非限制网络容量。

%% ========== 生物学启发 ==========

